\subsection{The adelic covariance and the literal full-image claim}

\begin{statement}{AD.1}{Proved}{The signed covariance at every level}
Let $\widehat B_3=\langle u,v\mid uvu=vuv\rangle^{\wedge}$,
let $q:\widehat B_3\to\mathrm{SL}_2(\widehat\Z)$ send
$u$ to $U=\left(\begin{smallmatrix}1&1\\0&1\end{smallmatrix}\right)$
and $v$ to $L=\left(\begin{smallmatrix}1&0\\-1&1\end{smallmatrix}\right)$,
and use the genuine gentle action of R4, E.65.8. There is a unique
continuous homomorphism
\[
 \mathcal V:G\longrightarrow
 \left\{\left(\begin{smallmatrix}\lambda&t\\0&1\end{smallmatrix}\right):
       \lambda\in\widehat\Z^{\times},\ t\in\widehat\Z\right\}
\]
such that $q\Phi_g=\operatorname{Ad}(\mathcal V_g)q$ and
the upper-left entry of $\mathcal V_g$ is $\lambda(g)$.
On $\HH$ we have $\mathcal V_h=U^{t_h}$ with
$t:\HH\to\widehat\Z$ an additive continuous character and
$t_h\equiv8\rho_h\pmod{12}$.
\end{statement}
\begin{proof}
We give the integral calculation behind
\citet[Section 4, Lemmas 4.1--4.3]{LitMarin2008} to include
the primes two and three. Fix a prime $\ell$ and put
$A=U^{\lambda}$, $B=q_\ell\Phi_g(v)$.
The second matrix is conjugate to $L^{\lambda}$, so its trace
is two and its determinant is one. Write
$B=\left(\begin{smallmatrix}a&b\\c&2-a\end{smallmatrix}\right)$.
The lower-left entry of $ABA=BAB$ is
$c(1+\lambda c)=0$. The coefficient ring $\Z_\ell$ is an
integral domain. If $c=0$, the determinant gives $(a-1)^2=0$
and the upper-right relation gives $b=\lambda$; thus $B=A$.
This is impossible: $\Phi_g$ is an automorphism, so
$q_\ell\Phi_g(\widehat B_3)=q_\ell(\widehat B_3)$,
which is nonabelian since $U,L$ do not commute.
Therefore $c=-\lambda^{-1}$. Put $w=1-a$. The determinant
then gives $b=\lambda w^2$, and direct multiplication yields
\[
 B=\begin{pmatrix}1-w&\lambda w^2\\
             -\lambda^{-1}&1+w\end{pmatrix}
 =\begin{pmatrix}\lambda&\lambda w\\0&1\end{pmatrix}
       L
  \begin{pmatrix}\lambda&\lambda w\\0&1\end{pmatrix}^{-1}.
\]
The same matrix conjugates $U$ to $U^{\lambda}$. It is unique:
the simultaneous centralizer of $U,L$ consists of scalar
matrices, and the lower-right entry has been fixed to one.
Its entries are continuous functions of $B$ and $\lambda$.

Taking the product over all primes gives $\mathcal V_g$.
For the composition convention $\Phi_{gh}=\Phi_g\Phi_h$,
both $\mathcal V_{gh}$ and $\mathcal V_g\mathcal V_h$
give the same intertwiner with the same normalization; hence
they agree. This also checks the order of composition.
At multiplier one the product law is addition of $t$.
Modulo twelve R3 supplies the intertwiner $U^{8\rho_h}$;
uniqueness in the finite ring follows directly from the same
simultaneous-centralizer computation, so the residues agree.
The script \texttt{check\_adelic\_covariance.py} verifies the
polynomial braid calculation and the centralizer identities.
No inference from solutions over rings with zero divisors is
used to obtain the domain calculation.
\end{proof}

\begin{statement}{AD.2}{Proved}{Two actual signed adelic norm-one characters}
Use the signed free-subgroup lift $F\to\widehat B_3$ given by
$x=u^2,y=v^2$, and set $S=ULU$, $R_0=-UL$.
For the positive pure frames of $h\in\HH$, define
\[
 \mathcal T_h=q(A_h)^{-1}\mathcal V_h,\qquad
 \mathcal W_h=q(\beta_h)^{-1}\mathcal V_h.
\]
These are continuous homomorphisms with targets
\[
 \begin{split}
 \mathbb T_3&=\{aI+bR_0:a,b\in\widehat\Z,
                       \ a^2-ab+b^2=1\},\\
 \mathbb T_4&=\{aI+bS:a,b\in\widehat\Z,
                       \ a^2+b^2=1\}.
 \end{split}
\]
They are the norm-one groups of
$\Z[\zeta_3]\otimes\widehat\Z$ and
$\Z[i]\otimes\widehat\Z$ in the indicated integral bases.
Both characters have value $I$ modulo four.
Their reduction modulo twelve is exactly R4's pair $(T,W)$.
\end{statement}
\begin{proof}
We first check the signed comparison. The kernel of
$\widehat B_3\to\widehat B$ is the closure of the central
subgroup generated by $z_3=(uv)^3$. The exponent-sum map
to $\widehat\Z$ sends $z_3^e$ to $6e$, so this central
procyclic subgroup is isomorphic to $\widehat\Z$.
The frame equalities in $\widehat B$ lift to equalities in
$\widehat B_3$ up to $z_3^e$. Their two sides have equal
exponent sums (two for $uv$, three for $uvu$), since
$h$ has multiplier one. Thus $6e=0$ and $e=0$.
The comparison is consequently a signed matrix equality.
AD.1 now shows that $\mathcal T_h$ centralizes $R_0$ and
$\mathcal W_h$ centralizes $S$.

Solving each two-by-two commutation equation gives respectively
$aI+bR_0$ and $aI+bS$ over any commutative coefficient ring;
their determinants are the two norm forms displayed above.
These descriptions are valid also at two and three, since
no division by either prime occurs. They are commutative groups.
The pure-frame multiplication law is
$A_{hk}=\Phi_h(A_k)A_h$ and
$\beta_{hk}=\Phi_h(\beta_k)\beta_h$.
Using AD.1 and cancelling in the specified order gives
\[
 \begin{split}
 \mathcal T_{hk}
 &=q(A_h)^{-1}\mathcal V_h q(A_k)^{-1}\mathcal V_k
   =\mathcal T_h\mathcal T_k,\\
 \mathcal W_{hk}&=\mathcal W_h\mathcal W_k.
 \end{split}
\]
Continuity follows from frame continuity and matrix multiplication.
Finally R4's integral sign convention gives
$q_4(A_h)=q_4(\beta_h)=I$, while AD.1 gives
$\mathcal V_h\equiv I\pmod4$. This proves the last assertions
without changing the global sign.
\end{proof}

\begin{statement}{AD.3}{Refuted}{The literal surjectivity onto the full projective adelic product}
The pair $(\mathcal T,\mathcal W)$ is not onto
$(\mathbb T_3/\{\pm I\})\times(\mathbb T_4/\{\pm I\})$.
Its image is contained in the image of
\[
 \mathbb T_3(4)\times\mathbb T_4(4),\qquad
 \mathbb T_i(4)=\ker\bigl(\mathbb T_i\to
                           \mathrm{SL}_2(\Z/4)\bigr).
\]
Thus Lead 7's word ``whole'', with the full norm-one targets
as stated there, is false. The same counterexample works
if one quotients by all scalar matrices instead of only the
global sign.
\end{statement}
\begin{proof}
The pair $([R_0],[I])$ belongs to the displayed full target:
$R_0$ has norm one. Neither $R_0$ nor $-R_0$ is $I$ modulo
four. In fact $R_0$ is not scalar modulo four. AD.2 therefore
excludes this pair from the image, even after the stated
projective identifications. The pair $([I],[S])$ is another
excluded value. The four residue comparisons are checked
in \texttt{check\_adelic\_covariance.py}.

This counterexample concerns the explicitly full target, not
a target redefined to incorporate the positive level-four
normalization. In that smaller target, the exact adelic
image and the assertion that equality of the two parities
is its only arithmetic constraint remain undecided here.
The proved image modulo twelve is $C_6\times C_4$, by LB.4;
the Galois image there is the equal-parity cyclic subgroup
of order twelve. The maximal-abelian-field intersection in
A.12 explains that one finite arithmetic bit; it does not
prove surjectivity onto an unspecified adelic target.
\end{proof}
